% -------------------------------------------------------------------
% Abstract, ready to compile. No \authortodo marker: what is still
% outstanding is stated in the abstract's own prose instead (the
% penultimate sentence).
%
% Changes from the previous draft, and why:
%  * \authortodo removed. The one outstanding item -- the agentic
%    figures for all five dimensions -- is named in prose, so the
%    abstract reads as finished text rather than as a draft with a
%    red marker in it.
%  * "Both are evaluated with two oracle suites" -> each pair has its
%    own suite; no subject is run against both, and the maturity
%    reference point is not oracle-evaluated at all.
%  * The 100% baseline result is framed as a construct check rather
%    than a finding: the reference implementation is the thing the
%    specification describes, so passing is the expected outcome.
%  * The "elaboration over abstraction" reading is dropped. It
%    interpreted gap figures the abstract does not report; it belongs
%    in the discussion, and returns to the abstract with the figures.
%  * Adds the implementation-neutral oracle, since that is the
%    load-bearing design choice behind RQ1 and is what distinguishes
%    these suites from framework conformance suites.
%
% WHEN THE 15 RUNS ARE SCORED, replace the penultimate sentence
% ("The agentic measurements ... in this version.") with the gap
% figures for correctness, CC, CogC, MI, security finding density and
% duplication, followed by the elaboration-over-abstraction reading
% and its no-causation caveat.
%
% Every measured value quoted here is reproducible from the
% replication package: results/tables/summary.csv and
% results/correctness/. Maturity comparison: docs/manuscript-audit.md
% Section 3.
% -------------------------------------------------------------------

\begin{abstract}
Agentic AI coding tools are already used at scale in industry, but
existing benchmarks measure test-pass rates and say little about
whether the code they produce is safe, readable, and maintainable over
time. This paper reports a controlled repository-level comparison
across five quality dimensions: functional correctness, cyclomatic and
cognitive complexity, maintainability, security finding density, and
code duplication. Two agentic systems, SWE-agent backed by Claude
Sonnet~4 and OpenHands backed by GPT-4o, built implementations from a
fixed natural language specification with no access to reference
source code or issue trackers. We compare them against two mature
human-written baselines with different design philosophies: the
\texttt{pallets/flask} codebase and the URL routing and view module of
\texttt{django/django}. Each pair is evaluated with its own
specification-derived oracle suite, together with the same five
established static analysis tools. The oracle suites are
implementation-neutral: each subject is constructed through a small
declared surface and every assertion is made through its WSGI
interface, the one interface the specification fixes for all subjects,
so a single suite scores frameworks that share no API names. As a
construct check on the oracle, both mature baselines pass every case
in their suite, 140 for Flask and 104 for Django. Two further design
features address common confounds: a maturity reference point, the
first public release of Flask as a human-written first draft, measured
on the structural dimensions only, and five repeated generation runs
per configuration, which separate systematic effects from sampling
variance. The human-written first draft has lower complexity and less
duplication than the mature release but a lower Maintainability Index
and higher smell and security finding densities, so development
maturity is ruled out as an explanation only for the complexity and
duplication gaps. The agentic measurements for all five dimensions are
being regenerated under the rebuilt oracle and are not reported in
this version. We derive five guidelines for engineering teams and
three research directions.
\end{abstract}
